%!TEX root = ../thesis.tex
%*******************************************************************************
%****************************** Chương 5 ***************************************
%*******************************************************************************

\chapter{THẢO LUẬN}
\label{chap:discussion}

\ifpdf
    \graphicspath{{Chapter5/Figs/Raster/}{Chapter5/Figs/PDF/}{Chapter5/Figs/}}
\else
    \graphicspath{{Chapter5/Figs/Vector/}{Chapter5/Figs/}}
\fi

Chương~\ref{chap:results} cho biết các mô hình đạt bao nhiêu. Chương này trả lời
vì sao, và lỗi còn lại nằm ở đâu. Mục~\ref{sec:dis-qlora} phân tích QLoRA sửa
được loại lỗi nào. Mục~\ref{sec:dis-ocr} xem xét cơ chế của QLoRA nhận biết OCR,
mục~\ref{sec:dis-remaining} và~\ref{sec:dis-audit} phân tích nguồn lỗi còn lại
bằng phân tích tự động và duyệt tay, mục~\ref{sec:dis-vivqa} xem xét lỗi trên
ViVQA. Mục~\ref{sec:dis-halluc} trả lời RQ4 về hallucination. Mục~\ref{sec:dis-v}
đối chiếu với bốn vấn đề tiếng Việt V1--V4, mục~\ref{sec:dis-examples} nêu ví dụ
định tính, mục~\ref{sec:dis-validity} bàn về giới hạn và tính hợp lệ, và
mục~\ref{sec:dis-ethics} bàn về khía cạnh đạo đức.

Mọi số trong chương chỉ dùng để mô tả, không dùng để chọn mô hình. Định nghĩa
kiểu lỗi và trạng thái OCR được ghi lại trước khi chạy phân tích
(mục~\ref{sec:met-error}).

%********************************** Mục 5.1 ************************************
\section{QLoRA sửa loại lỗi nào}
\label{sec:dis-qlora}

Bảng~\ref{tab:dis-buckets} phân loại các câu sai EM theo kiểu lỗi. Ở chế độ
zero-shot, lỗi thừa chữ (gold nằm trong câu trả lời) chiếm 25,7\% số lỗi trên
ViTextVQA và 21,3\% trên ViVQA. Sau QLoRA, tỉ lệ này giảm còn 13,6\% (B2) và 3,9\%
(C3). Như vậy phần lớn mức tăng của QLoRA đến từ việc mô hình học đúng định dạng
đáp án của bộ dữ liệu: ngắn, đúng cách viết. Kết luận này khớp với vấn đề V2 và
với phân tích zero-shot ở mục~\ref{sec:res-repro}. Sau QLoRA, lỗi còn lại dịch về
phía đọc sai ký tự (lệch vài ký tự, sai dấu) và sai nội dung.

\begin{table}[htbp]
\caption{Kiểu lỗi trên các câu sai EM (\% số câu sai của từng mô hình)}
\label{tab:dis-buckets}
\centering
\small
\begin{tabular}{p{0.21\textwidth}p{0.08\textwidth}p{0.08\textwidth}p{0.08\textwidth}p{0.08\textwidth}p{0.10\textwidth}p{0.07\textwidth}}
\toprule
\textbf{Mô hình} & \textbf{Tỉ lệ sai} & \textbf{Sai dấu} & \textbf{Thiếu chữ} & \textbf{Thừa chữ} & \textbf{Lệch $\leq$ 2 ký tự} & \textbf{Khác} \\
\midrule
ViTextVQA · B1 & 70,9 & 3,3 & 5,0 & \textbf{25,7} & 8,0 & 57,9 \\
ViTextVQA · B2 & 49,7 & 9,0 & 13,2 & 13,6 & 16,8 & 47,4 \\
ViTextVQA · B3-full & 46,6 & 9,6 & 15,5 & 14,5 & 16,5 & 43,9 \\
\midrule
ViVQA · B1 & 77,9 & 0,0 & 3,6 & \textbf{21,3} & 6,4 & 68,6 \\
ViVQA · C1 & 26,2 & 0,5 & 2,8 & 3,2 & 15,1 & 78,4 \\
ViVQA · C3 & 23,4 & 0,0 & 2,3 & 3,9 & 16,0 & 77,9 \\
\bottomrule
\end{tabular}
\end{table}

%********************************** Mục 5.2 ************************************
\section{Cơ chế của QLoRA nhận biết OCR}
\label{sec:dis-ocr}

Thí nghiệm $2\times2$ (Bảng~\ref{tab:res-2x2}) cho thấy OCR chỉ có ích khi mô hình
được huấn luyện cùng OCR. Một cách giải thích phù hợp: mô hình chưa thấy OCR khi
huấn luyện không biết dòng OCR đáng tin đến đâu, nên dùng nó không phân biệt và
sửa bao nhiêu câu thì cũng làm hỏng bấy nhiêu câu (16 và 16 trên dev300). Khi được
huấn luyện cùng OCR, adapter học được khi nào chép từ OCR và khi nào dựa vào ảnh.

Bảng~\ref{tab:dis-flips} chia các câu thay đổi từ B2 sang B3-full theo trạng thái
OCR. Hiệu số 705 câu được sửa trừ 395 câu bị làm hỏng bằng 310 câu, đúng bằng mức
tăng 3,09 điểm EM. Lợi ích tập trung ở các câu mà OCR chứa nguyên văn đáp án (sửa
326, hỏng 116). Ở các câu OCR không chứa đáp án, sửa và hỏng gần cân bằng (289 và
211).

\begin{table}[htbp]
\caption{Thay đổi từ B2 sang B3-full theo trạng thái OCR (toàn bộ test ViTextVQA)}
\label{tab:dis-flips}
\centering
\small
\begin{tabular}{p{0.34\textwidth}rrrrr}
\toprule
\textbf{Trạng thái OCR} & \textbf{n} & \textbf{Sửa được} & \textbf{Làm hỏng} & \textbf{Cùng đúng} & \textbf{Cùng sai} \\
\midrule
OCR chứa nguyên văn đáp án & 3.321 & \textbf{326} & 116 & 2.105 & 774 \\
OCR đúng chữ nhưng sai dấu & 1.149 & 90 & 68 & 506 & 485 \\
OCR có chữ, không chứa đáp án & 5.526 & 289 & 211 & 2.023 & 3.003 \\
OCR không đọc ra chữ & 32 & 0 & 0 & 11 & 21 \\
\midrule
Tất cả & 10.028 & 705 & 395 & 4.645 & 4.283 \\
\bottomrule
\end{tabular}
\end{table}

OCR cũng có thể làm mô hình phân tâm. Trong 138 câu (34,9\% số câu bị làm hỏng,
khoảng 1,4\% tập test), B2 trả lời đúng nhưng B3-full trả lời bằng một cụm khác
lấy từ dòng OCR. Đây là cái giá của việc đưa thêm chữ vào prompt, và là một lý do
prototype hiển thị chữ OCR để người dùng tự kiểm tra.

%********************************** Mục 5.3 ************************************
\section{Lỗi còn lại của B3-full: do OCR hay do mô hình}
\label{sec:dis-remaining}

Với mỗi câu ViTextVQA, nhóm xét đúng phần chữ OCR đã vào prompt. Chỉ 33,1\% câu có
đáp án gold nằm nguyên văn trong OCR; 11,5\% câu chỉ khớp khi bỏ dấu (OCR đọc đúng
chữ nhưng sai dấu); 55,1\% câu OCR có chữ nhưng không chứa nguyên văn đáp án; và
0,3\% câu OCR không đọc ra chữ nào. Đáng chú ý, EM của B3-full (53,35) cao hơn
nhiều so với tỉ lệ 33,1\%, nghĩa là ở nhiều câu mô hình tự đọc ảnh đúng mà OCR
không giúp được.

Bảng~\ref{tab:dis-source} xếp 4.678 câu sai của B3-full theo nguồn lỗi tự động.

\begin{table}[htbp]
\caption{Nguồn lỗi tự động của B3-full (4.678 câu sai)}
\label{tab:dis-source}
\centering
\small
\begin{tabular}{p{0.46\textwidth}rp{0.30\textwidth}}
\toprule
\raggedright \textbf{Nguồn lỗi} & \raggedright \textbf{\%} & \raggedright \textbf{Kiểu lỗi nhiều nhất} \tabularnewline
\midrule
\raggedright OCR không đưa nguyên văn đáp án vào prompt & \raggedright 81,0 & \raggedright khác 1.816; lệch ký tự 682; thiếu chữ 537 \tabularnewline
\raggedright OCR có đáp án, mô hình chép sai đoạn & \raggedright 10,5 & \raggedright thừa chữ 193; thiếu chữ 190 \tabularnewline
\raggedright OCR có đáp án, câu trả lời nằm ngoài OCR & \raggedright 8,5 & \raggedright khác 143; thừa chữ 123 \tabularnewline
\bottomrule
\end{tabular}
\end{table}

Con số 81\% cần được đọc thận trọng. Đây là cận trên của lỗi do OCR, vì nhóm ``OCR
không chứa đáp án'' gồm cả những câu cần suy luận (đáp án không được viết nguyên
văn trong ảnh) và những câu có gold nhiễu; mục~\ref{sec:dis-audit} tách các trường
hợp này bằng duyệt tay. Khi OCR có đáp án mà mô hình vẫn sai, lỗi chủ yếu là cắt
sai ranh giới cụm chữ: 383 trong 492 câu lấy thiếu hoặc thừa từ. Sai dấu từ OCR là
thách thức riêng của tiếng Việt: 393 trong 450 lỗi sai dấu của B3-full nằm ở
nhóm OCR không đưa nguyên văn đáp án. Theo loại câu hỏi, tỉ lệ lỗi do OCR cao nhất
ở loại liên hệ (90,9\%) và nội dung chữ (87,6\%), thấp nhất ở loại số và thời gian
(72,5\%).

%********************************** Mục 5.4 ************************************
\section{Duyệt tay 150 lỗi}
\label{sec:dis-audit}

Bảng~\ref{tab:dis-audit} tổng hợp kết quả duyệt tay mù 100 lỗi EM của B3-full và
50 lỗi EM của C3 (thiết kế ở mục~\ref{sec:met-error}). Cột cuối là tham chiếu từ
lần duyệt 200 lỗi của B2 ở giai đoạn trước, chỉ để mô tả: khác người gán, khác mô
hình và các khoảng tin cậy chồng nhau.

\begin{table}[htbp]
\caption{Kết quả duyệt tay lỗi EM (CI95 Wilson trong ngoặc)}
\label{tab:dis-audit}
\centering
\small
\begin{tabular}{p{0.32\textwidth}p{0.20\textwidth}p{0.20\textwidth}p{0.17\textwidth}}
\toprule
 & \textbf{B3-full (100 lỗi)} & \textbf{C3 (50 lỗi)} & \textbf{Tham chiếu B2 (ViTextVQA / ViVQA)} \\
\midrule
Lỗi EM thật ra đúng ý & \textbf{55,0} [45,2; 64,4] & \textbf{64,0} [50,1; 75,9] & 46 / 62 \\
Lỗi là hallucination & 5,0 [2,2; 11,2] & 16,0 [8,3; 28,5] & 9 / 5 \\
EM $\to$ độ đúng theo người đọc & 53,35 $\to$ \textbf{79,01} & 76,64 $\to$ \textbf{91,59} & \\
Hallucination trên mọi câu trả lời & 2,33\% & 3,74\% & \\
\bottomrule
\end{tabular}
\end{table}

\begin{table}[htbp]
\caption{Kiểu lỗi theo nhãn tay (\% số lỗi đã duyệt)}
\label{tab:dis-audit-types}
\centering
\small
\begin{tabular}{p{0.46\textwidth}rr}
\toprule
\textbf{Kiểu lỗi} & \textbf{B3-full} & \textbf{C3} \\
\midrule
Định dạng, độ dài (ý đúng) & \textbf{40,0} & --- \\
Đồng nghĩa, khác từ (ý đúng) & 4,0 & \textbf{36,0} \\
Đọc chữ sai & 29,0 & --- \\
Sai dấu thanh & 7,0 & --- \\
Sai vùng hoặc không có trong ảnh & 15,0 & 28,0 \\
Gold sai hoặc nhiễu & 5,0 & \textbf{28,0} \\
Đếm sai & --- & 6,0 \\
Khác & --- & 2,0 \\
\bottomrule
\end{tabular}
\end{table}

Có bốn nhận xét.

\paragraph{EM đánh giá thấp mô hình rất nhiều.} Hơn một nửa số lỗi EM của B3-full
thật ra đúng ý, nên độ đúng theo người đọc khoảng 79\% so với EM 53,35\%. Kiểu lỗi
lớn nhất là định dạng và độ dài (40\%): câu trả lời và gold cùng nội dung nhưng
khác cách viết, ví dụ ``combo 3'' so với gold ``3'', hoặc khoảng trắng quanh dấu
chấm trong địa chỉ email. Với C3, 36\% lỗi là đồng nghĩa, nên độ đúng theo người
đọc khoảng 92\%.

\paragraph{81\% lỗi ``do OCR'' không có nghĩa 81\% lỗi do OCR.} Trong 78 câu thuộc
nhóm OCR không chứa đáp án, chỉ 34 câu (43,6\%, CI95 [33,1; 54,6]) là lỗi đọc chữ
thật (đọc sai 27, sai dấu 7), còn 28 câu (35,9\%) là đúng ý. Tính trên cả mẫu, lỗi
đọc chữ thật chiếm 36\% (CI95 [27,3; 45,8]) số lỗi của B3-full, tức khoảng 17 điểm
EM. Đây là ước lượng thô cho phần mà một OCR tiếng Việt tốt hơn có thể sửa được.

\paragraph{Lỗi cắt sai ranh giới phần lớn được người đọc chấp nhận.} Trong 14 câu
mô hình chép sai đoạn từ OCR, 11 câu là định dạng hoặc độ dài với ý đúng.

\paragraph{Gold ViVQA nhiễu.} 28\% lỗi của C3 là do gold sai hoặc nhiễu, khớp với
các câu hỏi dịch máy khó hiểu ở mục~\ref{sec:dis-examples}.

Giới hạn của lần duyệt này: chỉ một người gán nên không tính được độ đồng thuận
$\kappa$ (bộ nhãn này đã đo $\kappa = 0{,}62$ giữa hai người ở lần duyệt trước);
50 câu ViVQA cho khoảng tin cậy rộng; và ``đúng ý'' là đánh giá chủ quan.

%********************************** Mục 5.5 ************************************
\section{Lỗi trên ViVQA}
\label{sec:dis-vivqa}

Với C3, lỗi phân bố theo loại câu hỏi như sau:

\begin{itemize}
  \item \textbf{Số lượng} (444 câu, sai 23,0\%): toàn bộ lỗi là đếm lệch, chủ yếu
  lệch một (bốn thành ba 11 câu, ba thành bốn 10 câu, một thành hai 10 câu).
  \item \textbf{Màu sắc} (625 câu, sai 15,7\%): nhầm màu gần nhau (nâu thành đen
  5 câu, đỏ thành cam 5 câu). Có 4 câu gold ``đỏ'' bị trả lời ``tía'', đúng lỗi
  dịch máy ``purple'' thành ``tía'' đã gặp khi duyệt dữ liệu.
  \item \textbf{Vật thể} (1.247 câu, sai 23,3\%): nhiều cặp đồng nghĩa hoặc khác
  mức chi tiết (máy tính và laptop, hải âu và con chim, tủ lạnh và tủ đá).
  \item \textbf{Vị trí} (685 câu, sai 30,7\%, khó nhất): lệch mức chi tiết (phòng
  và phòng ngủ, nhà vệ sinh và phòng tắm).
\end{itemize}

Phần lớn các lỗi vật thể và vị trí là đúng ý, như duyệt tay ở
mục~\ref{sec:dis-audit} xác nhận. Một thước đo chấp nhận đồng nghĩa sẽ phản ánh
tốt hơn chất lượng thật của mô hình trên ViVQA.

%********************************** Mục 5.6 ************************************
\section{Hallucination}
\label{sec:dis-halluc}

Bảng~\ref{tab:dis-pope} trình bày kết quả trên POPE-vi (thiết kế ở
mục~\ref{sec:met-pope}).

\begin{table}[htbp]
\caption{Kết quả POPE-vi (500 ảnh chưa thấy khi huấn luyện, 4.985 câu)}
\label{tab:dis-pope}
\centering
\small
\begin{tabular}{p{0.27\textwidth}p{0.08\textwidth}p{0.08\textwidth}p{0.08\textwidth}p{0.10\textwidth}p{0.10\textwidth}}
\toprule
\textbf{Cấu hình} & \textbf{F1 random} & \textbf{F1 pop.} & \textbf{F1 adv.} & \textbf{Tỉ lệ ``có''} & \textbf{Không đọc được} \\
\midrule
B1 · zero-shot, prompt thường & 84,60 & 83,17 & 81,19 & 40--43\% & 0,00\% \\
B1-ocr · zero-shot, prompt OCR & 83,39 & 81,99 & 80,59 & 38--41\% & 0,00\% \\
C3 & --- & --- & --- & --- & 91,43\% \\
\textbf{B3-full} & \textbf{88,89} & \textbf{85,65} & \textbf{83,06} & 48--55\% & 0,00\% \\
SH (adapter của prototype) & 77,70 & 72,83 & 71,17 & 77--89\% & 0,16\% \\
\bottomrule
\end{tabular}
\end{table}

\subsection{Kết luận theo luật}

Luật ghi trước so mỗi mô hình với mô hình nền zero-shot dùng cùng dạng prompt và
coi mô hình ``không tăng hallucination'' nếu cận dưới CI95 của $\Delta$F1 lớn hơn
$-2$ ở cả ba cách chọn.

\begin{itemize}
  \item \textbf{B3-full}: $\Delta$F1 $+5{,}50$ / $+3{,}65$ / $+2{,}47$ (random /
  popular / adversarial) so với B1-ocr: \emph{không tăng hallucination}.
  \item \textbf{C3}: 91,43\% câu trả lời không phải ``có'' hay ``không'', vượt
  ngưỡng hợp lệ 5\%: \emph{chưa kết luận được}.
  \item \textbf{SH}: $\Delta$F1 $-5{,}69$ / $-9{,}16$ / $-9{,}42$, cận trên CI95
  dưới $-2$: \emph{tăng hallucination}. SH không thuộc H4 vì là adapter của
  prototype, nhưng kết quả này là một kiểm tra an toàn bắt buộc trước khi dùng.
\end{itemize}

Với B3-full đạt và C3 không đo được, H4 được kết luận là ủng hộ một phần.

\subsection{Phân tích sau, ngoài luật}

\paragraph{F1 che mất việc B3-full báo sai nhiều hơn.} Tỉ lệ báo sai (câu ``không''
bị trả lời ``có'') của B1-ocr là 2,7\% / 5,7\% / 8,9\%, còn của B3-full là 8,7\% /
16,1\% / 22,4\%. F1 của B3-full tăng vì recall tăng (73,5\% lên 86,9\%), bù cho
precision giảm. Nếu đo hallucination vật thể theo nghĩa hẹp, tức ``nói có vật khi
vật không có'', thì B3-full bịa nhiều hơn zero-shot, khoảng 2,5 lần ở cách chọn
adversarial. Một phần là vì mô hình zero-shot vốn dè dặt (tỉ lệ trả lời ``có''
khoảng 40\%, dưới mức 50\% của thiết kế POPE). Đây là một giới hạn phương pháp:
F1, thước đo được chọn từ đầu theo cách báo cáo của LLaVA-1.5, không đủ để kết
luận về hallucination; cần báo kèm tỉ lệ báo sai và tỉ lệ trả lời ``có''.

\paragraph{C3 quên cách trả lời có/không.} ViVQA không có câu hỏi có/không. Sau
QLoRA hai epoch, C3 trả lời câu có/không bằng danh từ (``bức ảnh'' 525 lần, ``cái
ghế'' 434 lần), chỉ 372 trên 4.985 câu là ``không''. Đây là một dạng quên kỹ năng
làm theo chỉ dẫn khi huấn luyện trên một tập đáp án hẹp.

\paragraph{SH thiên lệch ``có''.} SH trả lời ``có'' ở 3.773 trên 4.985 câu, tỉ lệ
báo sai 55--79\%. Đếm trong dữ liệu cho thấy cả train ViTextVQA lẫn train ViVQA
không có đáp án ``có'' nào, nên giả thuyết ``học lệch từ dữ liệu có/không'' bị
bác; nguyên nhân chưa rõ. Nhóm sửa ở lúc suy luận bằng ngưỡng $\tau^* = 4{,}85$
trên hiệu log xác suất giữa ``có'' và ``không'' (mục~\ref{sec:met-proto}), chọn
trên một tập hiệu chỉnh tách biệt với POPE-vi. Sau khi sửa, tỉ lệ trả lời ``có''
của SH về 41--45\% và SH đạt luật R-YB, không cần huấn luyện lại. Tỉ lệ báo sai
của SH sau khi sửa vẫn cao hơn B1-ocr khoảng 2--4 điểm.

\paragraph{Chen tiếng Anh.} B3-full trả lời bằng chữ tiếng Anh ``no'' ở 881 trên
4.985 câu (17,7\%), và gần như luôn đúng (802 câu ở câu có đáp án ``không'').

\paragraph{Vật bị báo sai nhiều nhất} ở mọi cấu hình là cốc, rồi đến ghế, bát và
chai. Một phần có thể do chú thích COCO thiếu nhãn, một giới hạn đã ghi trước.

\paragraph{Đối chiếu với duyệt tay.} Hallucination chiếm 5,0\% lỗi EM của B3-full
và 16,0\% lỗi của C3 (Bảng~\ref{tab:dis-audit}), tức khoảng 2,3\% và 3,7\% mọi
câu trả lời. Hai cách đo cùng cho thấy hallucination không phải loại lỗi chính
của các mô hình đã thích nghi trên VQA tiếng Việt.

%********************************** Mục 5.7 ************************************
\section{Bốn vấn đề tiếng Việt}
\label{sec:dis-v}

Bảng~\ref{tab:dis-v} đối chiếu bốn vấn đề V1--V4 nêu ở Chương~\ref{chap:intro} với
số đo cụ thể của khóa luận.

\begin{table}[htbp]
\caption{Bốn vấn đề tiếng Việt và số đo tương ứng}
\label{tab:dis-v}
\centering
\small
\begin{tabular}{p{0.20\textwidth}p{0.70\textwidth}}
\toprule
\raggedright \textbf{Vấn đề} & \raggedright \textbf{Số đo} \tabularnewline
\midrule
V1 · Dấu thanh khi đọc chữ &
OCR giữ đúng dấu: PaddleOCR 27,0\% so với VietOCR 37,7\% ($p = 0{,}0008$); 9,0\%
lỗi của B2 chỉ sai dấu; 11,5\% câu có OCR đúng chữ nhưng sai dấu; zero-shot EM
22,4 với đáp án có dấu so với 39,4 với đáp án Latin không dấu. \\
V2 · Sai ngôn ngữ, sai định dạng &
Zero-shot trả lời dài: 2,9 so với 1,6 từ (ViVQA), 5,4 so với 3,7 từ (ViTextVQA);
lỗi thừa chữ giảm từ 25,7\% xuống 13,6\% số lỗi sau QLoRA; B3-full trả lời ``no''
tiếng Anh ở 17,7\% câu có/không. \\
V3 · Thứ tự chữ, ranh giới cụm &
Nhầm vùng chữ chiếm 23\% lỗi của B2; 383/492 lỗi khi OCR có đáp án là cắt sai
ranh giới cụm; F1 theo từ thấp hơn F1 theo âm tiết 1--3 điểm. \\
V4 · Hallucination &
2,3\% câu trả lời của B3-full; POPE-vi $\Delta$F1 từ $+2{,}5$ đến $+5{,}5$, nhưng
tỉ lệ báo sai tăng ở cách chọn adversarial. \\
\bottomrule
\end{tabular}
\end{table}

V2 được QLoRA sửa phần lớn. V1 và V3 là lý do chọn QLoRA nhận biết OCR, và cũng là
nơi còn nhiều lỗi nhất: khoảng 17 điểm EM của B3-full là lỗi đọc chữ thật. V4 được
đo riêng bằng POPE-vi và duyệt tay.

%********************************** Mục 5.8 ************************************
\section{Ví dụ định tính}
\label{sec:dis-examples}

Bảng~\ref{tab:dis-examples} nêu các ví dụ được rút tất định (seed 42) từ các nhóm
lỗi, minh họa những hiện tượng đã đo ở trên.

\begin{table}[htbp]
\caption{Ví dụ định tính trên test ViTextVQA và ViVQA}
\label{tab:dis-examples}
\centering
\footnotesize
\begin{tabular}{p{0.14\textwidth}p{0.20\textwidth}p{0.12\textwidth}p{0.12\textwidth}p{0.12\textwidth}p{0.10\textwidth}}
\toprule
\raggedright \textbf{Hiện tượng} & \raggedright \textbf{Câu hỏi} & \raggedright \textbf{Gold} & \raggedright \textbf{B2} & \raggedright \textbf{B3-full} & \raggedright \textbf{Ghi chú} \tabularnewline
\midrule
\raggedright OCR sửa đúng tên riêng & \raggedright trạm xăng này có tên là gì? & \raggedright mipec & \raggedright mipeo & \raggedright \textbf{mipec} & \raggedright chép đúng từng ký tự \tabularnewline
\raggedright OCR sửa đúng số & \raggedright bức hình được chụp vào năm nào? & \raggedright 2016 & \raggedright 2017 & \raggedright \textbf{2016} & \raggedright số dễ đọc nhầm \tabularnewline
\raggedright OCR sửa cả chữ lẫn độ dài & \raggedright trường thpt này mang tên gọi là gì? & \raggedright đào duy từ & \raggedright thành phố đào duy trì & \raggedright \textbf{đào duy từ} & \tabularnewline
\raggedright OCR làm phân tâm & \raggedright \ldots\ liên hệ với người nào? & \raggedright cô tuyết & \raggedright \textbf{cô tuyết} & \raggedright có tuyết & \raggedright OCR sai dấu \tabularnewline
\raggedright Cắt sai ranh giới & \raggedright món cơm gà hải nam làm theo nước nào? & \raggedright kiểu singapore & \raggedright \textbf{kiểu singapore} & \raggedright singapore & \raggedright đúng ý \tabularnewline
\raggedright Sai dấu thanh & \raggedright di tích lịch sử này tên gì? & \raggedright nhà tù hòa bình & \raggedright nhà từ hòa bình & \raggedright nhà tu hòa bình & \raggedright cả hai sai \tabularnewline
\raggedright Định dạng & \raggedright máy mài góc \ldots\ thuộc combo nào? & \raggedright 3 & \raggedright combo 3 & \raggedright combo 3 & \raggedright đúng ý \tabularnewline
\raggedright Chuỗi dài & \raggedright địa chỉ email của công ty là gì? & \raggedright tnhhanphat\-vn@gmail . com & \raggedright tnhanhphat\-vn@\ldots & \raggedright tnhphanhat\-vn@\ldots & \raggedright lệch vài ký tự \tabularnewline
\bottomrule
\end{tabular}
\end{table}

Trên ViVQA, câu hỏi ``đồng hồ và tranh cũ dòng những gì'' với gold ``tường'' là
một ví dụ câu hỏi dịch máy khó hiểu: cả B1 (``đồng hồ và tranh cũ dòng những bức
tranh'') lẫn C3 (``con chim'') đều trả lời sai. Câu hỏi ``màu của nho là gì'' với
gold ``màu đỏ'' được C3 trả lời ``màu tía''; ranh giới giữa đỏ, tía và đen của
nho vốn mơ hồ.

%********************************** Mục 5.9 ************************************
\section{Giới hạn và tính hợp lệ}
\label{sec:dis-validity}

\paragraph{Một seed cho mô hình chính.} B3-full và C3 chỉ chạy với seed 42 vì mỗi
lần huấn luyện B3-full mất gần 13 giờ GPU. Độ ổn định theo seed chỉ được kiểm trên
tập con 8.000 câu (ba seed, $+3{,}57 \pm 0{,}76$ EM). Với $p = 6{,}7\times10^{-21}$
trên 10.028 câu, khả năng kết luận B3-full hơn B2 đảo chiều theo seed là thấp,
nhưng độ lớn chính xác của mức tăng có thể dao động.

\paragraph{EM quá khắt khe.} Duyệt tay cho thấy 55\% lỗi EM của B3-full và 64\%
lỗi của C3 thật ra đúng ý. EM vì vậy đánh giá thấp chất lượng thật; nhóm bổ sung
ANLS và duyệt tay, nhưng chưa có đánh giá người quy mô lớn.

\paragraph{So sánh với kết quả công bố.} Bài báo ViTextVQA không công bố hàm chuẩn
hóa đáp án. Mức chênh với ViTextBLIP-2 lớn đến mức khó do chuẩn hóa gây ra, nhưng
con số tuyệt đối chỉ mang tính tham khảo.

\paragraph{H2 về ANLS chưa đạt ngưỡng.} OCR giúp trả lời đúng nguyên chuỗi (EM) rõ
hơn là giúp gần đúng theo ký tự (ANLS $+1{,}87$ so với ngưỡng 3). Ngưỡng được đặt
trước khi có số liệu và nhóm giữ nguyên.

\paragraph{Phụ thuộc OCR.} B3 phụ thuộc vào OCR: bỏ OCR lúc suy luận thì về mức
không OCR. Chất lượng OCR tiếng Việt cho chữ cảnh là nút thắt chính, và lỗi
EXIF (mục~\ref{sec:res-exif}) cho thấy đường OCR cần được kiểm chặt.

\paragraph{Dữ liệu.} Câu hỏi ViVQA được dịch máy và một phần gold nhiễu; khoảng một
nửa ảnh test ViVQA có trong train. Nhóm đo và báo cáo các vấn đề này, đồng thời
kiểm trên tập sạch, nhưng không sửa dữ liệu để giữ khả năng so sánh. Bộ phân loại
câu hỏi \texttt{qt-v1} có độ chính xác 80,7\%, nên bảng theo loại câu hỏi chỉ mang
tính mô tả.

\paragraph{POPE-vi.} Chú thích COCO có thể thiếu nhãn vật thể, và mô hình nền có
thể đã thấy ảnh COCO khi tiền huấn luyện. POPE-vi chỉ đo hallucination vật thể
trên ảnh đời thường, chưa đo hallucination chữ trên ảnh có chữ.

\paragraph{Người gán nhãn.} Lần duyệt tay cuối chỉ có một người gán.

\paragraph{Độ mạnh kiểm định trên test2000.} Các ablation trên 2.000 câu chỉ phát
hiện được chênh lệch từ khoảng 1,7 điểm EM; vì vậy kết luận ``không khác biệt'' ở
ablation cấu hình LoRA là kết luận yếu.

%********************************** Mục 5.10 ***********************************
\section{Khía cạnh đạo đức và sử dụng có trách nhiệm}
\label{sec:dis-ethics}

\paragraph{Quyền riêng tư.} Ảnh người dùng tải lên có thể chứa khuôn mặt, giấy tờ
hoặc thông tin cá nhân. Prototype không lưu ảnh hay câu hỏi, chỉ ghi thời gian xử
lý.

\paragraph{Thiên lệch ngôn ngữ.} Mô hình nền được tiền huấn luyện thiên về tiếng
Anh. Khóa luận đo được thiên lệch này (chênh 17 điểm EM giữa đáp án có dấu và
không dấu ở zero-shot; mô hình chen tiếng Anh) và giảm nó bằng thích nghi, nhưng
không loại bỏ hoàn toàn.

\paragraph{Hallucination và độ tin cậy.} Câu trả lời trôi chảy nhưng sai là rủi ro
chính. Prototype hiển thị chữ OCR và mức tin cậy để người dùng tự kiểm tra, và
chỉ hiện mức (Cao, Trung bình, Thấp) vì độ tin cậy chưa được hiệu chỉnh tốt
(ECE 0,198). Hệ thống không nên được dùng cho các tình huống nhạy cảm như y tế
hoặc pháp lý khi chưa có người kiểm tra.

\paragraph{Giấy phép dữ liệu.} ViTextVQA được dùng qua kho có kiểm soát truy cập,
theo giấy phép CC BY-NC 3.0 chỉ cho mục đích nghiên cứu phi thương mại; tập test
có đáp án không được phát tán lại.
